\chapter{Services REST}\label{chap:rest}

\section{Conventions communes}
\begin{itemize}
  \item \textbf{Ressources} nommées au pluriel (\code{/candidatures}, \code{/stages}), sous-ressources pour les
        compositions (\code{/stages/\{id\}/evaluations}), filtres en paramètres de requête (\code{?matricule=\&statut=}).
  \item \textbf{Méthodes HTTP} : \code{GET} lecture (sûre, idempotente), \code{POST} création, \code{PATCH} modification
        partielle (statut, complément), \code{DELETE} annulation.
  \item \textbf{Format} : JSON UTF-8, dates ISO-8601, champs nuls omis.
  \item \textbf{Création} : \code{201 Created} avec l'en-tête \code{Location} de la ressource créée.
  \item \textbf{Authentification} par clé d'API (en-tête \code{X-API-Key}) pour les applications et scripts, ou
        HTTP Basic pour l'interface web ; le périmètre est restreint automatiquement selon l'acteur
        (un étudiant ne voit que ses candidatures, une entreprise que celles qui la concernent).
\end{itemize}

\begin{table}[htbp]
\centering
\caption{Codes HTTP retournés}\label{tab:codes}
\begin{tabularx}{\textwidth}{@{}>{\raggedright\arraybackslash}p{3.3cm}>{\raggedright\arraybackslash}p{5.4cm}Y@{}}
\toprule
\textbf{Code} & \textbf{Erreur} & \textbf{Cas} \\
\midrule
200 OK & — & Lecture, modification, annulation réussies \\
201 Created & — & Création (candidature, notification, stage, compte rendu, évaluation, instance de processus) \\
400 Bad Request & \code{DONNEES\_INVALIDES}, \code{JSON\_INVALIDE} & Champs obligatoires manquants, format ou bornes non respectés, JSON mal formé \\
401 Unauthorized & \code{NON\_AUTHENTIFIE}, \code{CLE\_API\_INVALIDE} & Identifiants ou clé d'API absents, invalides ou révoqués \\
403 Forbidden & \code{ACCES\_REFUSE} & Rôle non autorisé ou ressource hors périmètre \\
404 Not Found & \code{RESSOURCE\_INTROUVABLE} & Identifiant inconnu \\
409 Conflict & \code{CONFLIT}, \code{TRANSITION\_INTERDITE}, \code{DECISION\_VIA\_\allowbreak PROCESSUS} & Doublon, transition d'état interdite, décision contournant le processus \\
422 Unprocessable Entity & \code{REGLE\_METIER} & Règle de gestion (annulation d'une candidature déjà traitée…) \\
500 / 502 & \code{ERREUR\_INTERNE}, \code{ERREUR\_SERVICE\_AMONT} & Erreur inattendue, faute d'un service SOAP appelé \\
503 Service Unavailable & \code{SERVICE\_INDISPONIBLE} & Service indisponible (en-tête \code{Retry-After}) \\
\bottomrule
\end{tabularx}
\end{table}

Toutes les erreurs ont le même corps JSON, produit par un \code{ExceptionMapper} unique :
\begin{lstlisting}[language=json, caption={Corps d'erreur uniforme (400)}]
{
  "status": 400,
  "erreur": "DONNEES_INVALIDES",
  "message": "Données invalides",
  "details": [ "Le champ 'matricule' est obligatoire", "Le champ 'offreId' est obligatoire" ],
  "horodatage": "2026-10-23T18:22:16.985909069Z"
}
\end{lstlisting}

\section{Service 4 — Gestion des candidatures}
\textbf{Base} : \code{/api/candidatures}

\begin{table}[htbp]
\centering\small
\caption{Points d'accès du service Candidatures}
\begin{tabularx}{\textwidth}{@{}llYl@{}}
\toprule
\textbf{Méthode} & \textbf{URI} & \textbf{Description} & \textbf{Codes} \\
\midrule
POST & \code{/candidatures} & Dépôt (enregistrement) d'une candidature ; refus des doublons et au-delà de 5 candidatures en cours. Réservé à l'orchestrateur & 201, 400, 409, 422 \\
GET & \code{/candidatures} & Recherche (\code{matricule}, \code{offreId}, \code{entrepriseId}, \code{statut}) & 200 \\
GET & \code{/candidatures/etudiant/\{matricule\}} & Candidatures d'un étudiant & 200, 403 \\
GET & \code{/candidatures/\{id\}} & Détail & 200, 403, 404 \\
GET & \code{/candidatures/\{id\}/historique} & Historique horodaté des changements de statut & 200, 404 \\
PATCH & \code{/candidatures/\{id\}} & Complément du dossier (lettre, CV) par l'étudiant & 200, 422 \\
PATCH & \code{/candidatures/\{id\}/statut} & Changement de statut \code{\{statut, motif\}} contrôlé par l'automate & 200, 400, 409 \\
DELETE & \code{/candidatures/\{id\}} & Annulation (statut \code{ANNULEE}), uniquement tant que la candidature est \code{DEPOSEE} & 200, 403, 422 \\
\bottomrule
\end{tabularx}
\end{table}

\begin{figure}[htbp]
\centering
\begin{tikzpicture}[font=\sffamily\small, node distance=12mm and 28mm, >={Latex[length=2mm]},
  st/.style={draw=black!45, rounded corners=8pt, fill=#1, minimum width=3.2cm, minimum height=8mm}]
  \node[st=soap] (dep) {DEPOSEE};
  \node[st=bpmn, right=of dep] (eval) {EN\_COURS\_EVALUATION};
  \node[st=rest, right=22mm of eval, yshift=8mm] (acc) {ACCEPTEE};
  \node[st=black!6, below=14mm of eval, xshift=-24mm] (ann) {ANNULEE};
  \node[st=black!6, right=22mm of eval, yshift=-8mm] (ref) {REFUSEE};
  \node[left=10mm of dep] (init) {};
  \draw[->, thick] (init) -- node[above, font=\scriptsize]{POST} (dep);
  \draw[->, thick] (dep) -- node[above, font=\scriptsize]{prise en charge} (eval);
  \draw[->, thick] (eval) -- node[above, font=\scriptsize, sloped]{décision} (acc);
  \draw[->, thick] (eval) -- node[below, font=\scriptsize, sloped]{décision} (ref);
  \draw[->, thick] (dep) |- node[pos=0.25, left, font=\scriptsize]{DELETE} (ann);
  \draw[->, thick] (dep.north) to[bend left=22] node[above, font=\scriptsize]{refus direct} (ref.north);
  \draw[->, thick, black!50] (eval) edge[loop above] node[font=\scriptsize]{complément (PATCH)} (eval);
\end{tikzpicture}
\caption{Automate des statuts d'une candidature (états finals en gris et vert)}\label{fig:automate}
\end{figure}

\begin{lstlisting}[language=json, caption={Ressource candidature (GET \code{/api/candidatures/1})}]
{
  "id": 1,
  "matricule": "E2026001",
  "nomEtudiant": "Yasmine Benali",
  "offreId": 1,
  "titreOffre": "Assistant chef de projet AMOA",
  "entrepriseId": 1,
  "lettreMotivation": "Très motivée par la modélisation des processus métier.",
  "cvUrl": "https://cv.example/E2026001.pdf",
  "scoreCompatibilite": 100,
  "statut": "ACCEPTEE",
  "motif": "Bienvenue !",
  "processInstanceId": "5",
  "dateDepot": "2026-10-23T18:18:00.699389",
  "dateMiseAJour": "2026-10-23T18:18:03.470788"
}
\end{lstlisting}

\section{Service 5 — Notifications}
\textbf{Base} : \code{/api/notifications}

\begin{table}[htbp]
\centering\small
\caption{Points d'accès du service Notifications}
\begin{tabularx}{\textwidth}{@{}llYl@{}}
\toprule
\textbf{Méthode} & \textbf{URI} & \textbf{Description} & \textbf{Codes} \\
\midrule
POST & \code{/notifications} & Envoi d'une notification typée \code{\{type, destinataireType, destinataireId, email, candidatureId, parametres\}} ; le gabarit est rendu côté serveur & 201, 400 \\
GET & \code{/notifications} & Boîte de réception (filtres \code{destinataireType}, \code{destinataireId}, \code{candidatureId}, \code{type}) & 200 \\
GET & \code{/notifications/\{id\}} & Détail (destinataire uniquement) & 200, 403, 404 \\
PATCH & \code{/notifications/\{id\}/lue} & Marquer comme lue & 200 \\
POST & \code{/notifications/rappels} & Campagne de rappels des dates limites \code{\{joursAvant\}} & 200, 400 \\
\bottomrule
\end{tabularx}
\end{table}

Les messages sont produits à partir de gabarits paramétrés (\code{\{nomEtudiant\}}, \code{\{titreOffre\}},
\code{\{motif\}}…). Les types couvrent les besoins du sujet : confirmation de dépôt, information de l'entreprise,
acceptation, refus, demande de complément, rappels de date limite, ainsi que le rejet automatique, la relance de
l'entreprise et l'information en cas d'incident.

Le canal e-mail envoie de \textbf{vrais messages} en SMTP (Jakarta Mail), au format texte et HTML aux couleurs de
l'INPT. En démonstration, le serveur SMTP est \textbf{Mailpit} : il capture tous les messages dans une boîte de
réception consultable sur \code{http://localhost:8025}, sans qu'aucun e-mail ne sorte de la machine
(figure~\ref{fig:mailpit}). Chaque notification reste archivée en base avec son statut d'envoi (\code{ENVOYEE},
\code{ECHEC\_ENVOI} si le serveur SMTP ne répond pas, \code{JOURNALISEE} si l'envoi est désactivé), ce qui la rend
aussi consultable dans l'interface. Pour un vrai serveur, il suffit de changer \code{mail.smtp.hote}.

La \textbf{campagne de rappels} illustre la composition de services : le service REST interroge le service SOAP Offres
(\code{listerOffresExpirantBientot}), le service SOAP Étudiants (\code{listerEtudiants(statut=ACTIF)}) puis
\code{verifierCompatibilite} pour chaque couple, et n'envoie qu'un rappel par étudiant et par offre. Elle s'exécute
automatiquement toutes les 24 heures et peut être déclenchée à la demande.

\section{Service 6 — Suivi des stages}
\textbf{Base} : \code{/api/stages}

\begin{table}[htbp]
\centering\small
\caption{Points d'accès du service Suivi des stages}
\begin{tabularx}{\textwidth}{@{}llYl@{}}
\toprule
\textbf{Méthode} & \textbf{URI} & \textbf{Description} & \textbf{Codes} \\
\midrule
POST & \code{/stages} & Enregistrement d'un stage accepté (un seul par candidature) & 201, 400, 409 \\
GET & \code{/stages} & Recherche (\code{matricule}, \code{entrepriseId}, \code{statut}, \code{candidatureId}) & 200 \\
GET & \code{/stages/\{id\}} & Détail & 200, 403, 404 \\
GET & \code{/stages/\{id\}/statut} & Statut, dates, durée, jours restants et progression calculée & 200 \\
PATCH & \code{/stages/\{id\}/statut} & Transition \code{PLANIFIE→EN\_COURS→TERMINE}, ou \code{INTERROMPU} & 200, 409 \\
POST & \code{/stages/\{id\}/comptes-rendus} & Dépôt d'un compte rendu périodique & 201, 400, 422 \\
GET & \code{/stages/\{id\}/comptes-rendus} & Liste des comptes rendus & 200 \\
POST & \code{/stages/\{id\}/evaluations} & Évaluation (\code{MI\_PARCOURS} ou \code{FINALE}, note sur 20) & 201, 400, 409, 422 \\
GET & \code{/stages/\{id\}/evaluations} & Liste des évaluations & 200 \\
GET & \code{/stages/etudiant/\{matricule\}/historique} & Stages réalisés, nombre de comptes rendus, évaluations et moyenne & 200, 403 \\
\bottomrule
\end{tabularx}
\end{table}

\begin{lstlisting}[language=json, caption={GET \code{/api/stages/1/statut}}]
{ "stageId": 1, "statut": "TERMINE", "dateDebut": "2027-02-01", "dateFin": "2027-07-31",
  "dureeJours": 180, "joursRestants": 281, "progression": 100 }
\end{lstlisting}

L'évaluation finale clôture le stage (\code{TERMINE}) ; à la fin ou à l'interruption d'un stage, le service appelle
\code{changerStatutEtudiant(ACTIF)} du service SOAP Étudiants pour rendre l'étudiant à nouveau éligible.

\section{API d'orchestration}
\textbf{Base} : \code{/api/processus}. Elle expose le processus BPMN aux acteurs humains.

\begin{table}[htbp]
\centering\small
\caption{Points d'accès de l'orchestration}
\begin{tabularx}{\textwidth}{@{}llYl@{}}
\toprule
\textbf{Méthode} & \textbf{URI} & \textbf{Description} & \textbf{Rôles} \\
\midrule
POST & \code{/processus/candidatures} & Soumission : démarre une instance ; renvoie son état (rejet motivé, attente de décision ou incident) & Étud., Resp. \\
GET & \code{/processus/taches} & Tâches humaines en attente (filtrées par entreprise) & Ent., Resp. \\
POST & \code{/processus/taches/\{id\}/prise-en-charge} & L'entreprise commence l'examen : tâche attribuée, candidature \code{EN\_COURS\_EVALUATION} & Ent., Resp. \\
POST & \code{/processus/taches/\{id\}/decision} & Décision \code{ACCEPTEE} / \code{REFUSEE} / \code{COMPLEMENT} + commentaire, dates, tuteur & Ent., Resp. \\
POST & \code{/processus/taches/\{id\}/incident} & Traitement d'un incident : \code{RELANCER} ou \code{ABANDONNER} & Resp. \\
GET & \code{/processus/instances} & Instances récentes & Resp. \\
GET & \code{/processus/instances/\{id\}} & État, variables, tâches, activités actives et trace d'exécution & Concernés \\
GET & \code{/processus/definition} & Modèle BPMN 2.0 (XML) & Tous \\
\bottomrule
\end{tabularx}
\end{table}
